\chapter*{前言}
\addcontentsline{toc}{chapter}{前言}

\section*{第一张地图：我们究竟在研究什么}

我希望先和你约定一件事：在翻开这本书以后，我们暂时不把“智能”当作某个模型在某张榜单上的分数。分数当然有用，但它只描述一次被截取的表现；它没有告诉我们，这个系统昨天经历了什么，今天为何仍是同一个主体，明天又会因今天的选择发生怎样的改变。若一个人工系统每次启动都从近乎相同的参数出发，在会话结束后便失去自己的经历，那么它更接近一件能力很强的仪器，而不是一个持续存在的智能体。

本书研究的是后一个问题。我们要问：一个人工主体怎样把此刻的激活保存成经历，把反复出现的经历压缩成结构，把结构变成技能与身体习惯，又怎样让这些变化不至于破坏身份连续？它怎样知道自己正处在过载、冲突或失控边缘？怎样把稳定能力交给较快、较局部的控制回路，同时把意外残差送回较慢、较全局的认知层？当它拥有文件、工具、机器人身体和其他主体时，“它自己”的边界又画在哪里？

这些问题看起来分散，其实由同一个闭环贯穿：
\[
\text{世界}\longrightarrow\text{有限观测}\longrightarrow\text{内部结构}
\longrightarrow\text{行动}\longrightarrow\text{新世界}\longrightarrow\text{残差}.
\]
智能并不站在世界之外。它只能取得有限观测，只能形成有偏差的内部结构，也只能通过有延迟、有风险的身体改变现实。现实返回的残差，决定哪些猜想应当保留、哪些技能应当重练、哪些自我解释必须放弃。我们把这条原则称为现实权限：预测可以先行，行动可以尝试，但观察到的现实不能被预测结果悄悄覆盖。

你会在书中遇到许多新名称：CSO、Agent-CSO、三相记忆、SPC、AHC、TCE、SGC、FCS、BPCC、KRA。请不要把它们当作待背诵的部件表。每个名称都只是一张临时标签，标签下面必须有状态、输入、输出、时间尺度和失败条件。例如，SPC 若被称为自我感知压缩器，我们就必须追问：它压缩什么分布？估计误差如何表示？观察器延迟怎样进入控制稳定性？如果答不出来，这个术语便没有增加知识。

本书的第一条写作纪律因此是：先把直觉放进动力学，再允许它成为概念。我们会借助控制论、信息论、快慢系统、持续学习、具身认知、模型预测控制与混杂系统。这些已有理论构成地基，但不替本书的新构造背书。CSO 等概念仍是待验证的综合框架；“气态记忆”“认知地貌”“引力中心”等说法首先是帮助思考的比喻，绝不能因为像物理学语言便被误认为物理定律。

\begin{narration}
读到这里，我们得到第一张地图：研究对象不是孤立模型，而是一个在有限观察、有限容量、有限身体和真实反馈中延续的生命周期过程。
\end{narration}

\clearpage

\section*{第二张地图：怎样阅读一部尚在形成中的理论}

这本书来自一场递进式讨论。《比较智能差异.md》不是一篇预先写好的论文，而是一条不断修正自身的思考轨迹：先讨论不同尺度智能的差异，再讨论智能空间、语义类和生成基；随后把注意力转向三相记忆、自我、认知控制与 AgentOS；最后又用 CSO 与 Agent-CSO 重新整理早期“复杂度迁移”的说法，并把身体运行时、SGC/FCS、BPCC/KRA 和生命周期工程接入同一框架。因此，较晚形成的定义有时会修正较早的直觉。

我们没有把这种修正史抹平。相反，我会在关键位置告诉你：哪些概念被保留，哪些只是宏观近似，哪些已经被更精确的本体取代。例如，“复杂度迁移”仍能描述运行负担从显式认知下沉到技能或工具的宏观现象，但复杂度不是在模块间流动的物质。更精确的说法是：某个认知结构改变了表示、功能位置、时间尺度或承载边界。于是全书的顺序必须是先定义 CSO，再定义 Agent-CSO，最后才谈 CSO 迁移和拓扑沉降。

阅读时可以不断做四个检验。第一是对象检验：公式中的每个符号究竟代表状态、轨迹、生成器，还是仅仅一个名称？第二是尺度检验：两个变量若更新速度相差几个数量级，能否仍按同一个同步循环处理？第三是权限检验：预测、记忆、控制器和现实观察发生冲突时，谁有权覆盖谁？第四是反事实检验：若删除某一机制，系统会出现什么可观察差异？这四个问题能把漂亮的系统图还原成可实验的工程主张。

全书的概念依赖也按这四个检验安排。第一编只问智能是什么样的时空结构，并建立 CSO、承载映射和功能拓扑；第二编在拓扑暂时固定的条件下，引入三相记忆、自我慢变量和认知控制；第三编才让主体获得数字或物理身体；第四编把理论拆成可实施路线；后面的控制、培育和心理工程，则研究一个持续主体怎样稳定、成长以及从异常中恢复。终编把同一语言谨慎扩展到组织与文明，但不会假设组织必然拥有单一自我，更不会由功能自反推出现象意识。

我还要说明这本书刻意不做什么。它不提供一套可直接复制的产品代码，不把人类心理诊断机械地移植给机器，也不宣称已经解决人工意识。它给出的是一组结构假说、数学骨架、接口原则和工程判据。真正的理论成熟，必须依赖可重复实验：长期记忆是否改善跨月任务而不放大偏见？慢变量保护是否提高身份连续性却阻碍必要适应？局部能力下沉是否降低显式认知负担而不隐藏危险残差？这些问题都应当有失败答案的可能。

如果你是一位系统工程师，可以把本书看成持续 Agent 的架构推演；如果你研究控制或学习理论，可以把它看成一个等待形式化的多时间尺度混杂系统；如果你关心智能的认识论，那么请始终记住：内部世界不是现实副本，而是主体在有限生命史中形成的、可行动但可错的结构。

\begin{narration}
第二张地图告诉我们怎样前进：尊重概念的形成顺序，区分已有理论与新假说，让现实残差保留否决权，并把每个隐喻最终兑换成可测量的关系。
\end{narration}
